\section{Base de datos}

Motor PostgreSQL. La base de trabajo es \codigo{mpaz} en Neon; el servidor
del colegio usa PostgreSQL 16 local con base \codigo{mpaz\_matricula}. El
portal usa doble camino de conexión (endpoint del servicio web y conexión
directa) con comprobación \codigo{SELECT 1} al arrancar
(\codigo{ConexionBaseDatos.cs}).

Todos los objetos del portal viven en el esquema \codigo{mpaz\_matricula}
con \codigo{Search Path=mpaz\_matricula,public}; la extensión
\codigo{pg\_trgm} es compartida en \codigo{public}. El esquema
\codigo{mpaz\_auth} concentra Identity y OpenIddict; \codigo{mpaz\_asistencia}
reserva el módulo de asistencia (MAT-15, sin flujo de datos aún).

\subsection{Instalación}

Instalación nueva en tres partes, en orden, desde
\codigo{database/modulos/matricula/instalacion/}:
\codigo{01-estructura-bd.sql} $\rightarrow$ \codigo{02-programacion-pl.sql}
$\rightarrow$ \codigo{03-datos-iniciales.sql}. Para bases existentes se usa
el aplicador \codigo{database/aplicar.sh} con los pasos numerados de
\codigo{comun/pasos/} y \codigo{modulos/matricula/pasos/}.

\subsection{Pasos del esquema de matrícula}

\begin{longtable}{@{}p{0.14\textwidth}p{0.78\textwidth}@{}}
\toprule
\textbf{Paso} & \textbf{Contenido}\\
\midrule
\codigo{01} & DDL institucional: tablas, restricciones, triggers y vistas\\
\codigo{03} & Matriz inicial de permisos por rol\\
\codigo{04} & Privilegios mínimos del rol \codigo{mpaz\_app}\\
\codigo{05} & Año escolar configurable (\codigo{activo\_matricula})\\
\codigo{06} & Validación de matrícula\\
\codigo{07} & \codigo{search\_path} en funciones\\
\codigo{08} & Certificados emitidos\\
\codigo{09} & Auditoría de \codigo{funcionario}\\
\codigo{10--11} & Cursos de parvularia y letras A/B (pre-kínder a 8º básico)\\
\codigo{12--13} & Vista de alertas de vencimiento y tabla \codigo{alerta\_revision}\\
\codigo{14} & Roles SECRETARIA/DIRECCION/INSPECTORIA/PIE (Johan 18-11-2026)\\
\codigo{15} & Roles UTP/CONVIVENCIA y TENS completo\\
\bottomrule
\end{longtable}

\subsection{Grupos de tablas}

\begin{itemize}[leftmargin=*,itemsep=0.2em]
\item \textbf{Identificación}: \codigo{estudiante}, persona y familia
  (\codigo{estudiante\_persona}), apoderados y vínculos.
\item \textbf{Matrícula}: \codigo{matricula\_anual} (curso y año por fila),
  \codigo{ficha\_version} con \codigo{datos\_snapshot} JSONB y folio.
\item \textbf{Expediente}: \codigo{documento}, requisitos documentales,
  ubicación del original y estado de validación.
\item \textbf{Autorizaciones}: imagen, convivencia y medicación, con
  estados OTORGADA/RECHAZADA/REVOCADA/VENCIDA/REEMPLAZADA.
\item \textbf{Salud y medicación}: prescripciones y dosis (MAT-05, bloqueado
  por protocolo TENS).
\item \textbf{Seguridad y auditoría}: \codigo{funcionario}, \codigo{rol},
  \codigo{rol\_permiso}, \codigo{historial\_cambios} y
  \codigo{registros\_acceso}.
\item \textbf{Operación}: \codigo{anios\_escolares}, certificados,
  \codigo{alerta\_revision} y tablas de Identity en \codigo{mpaz\_auth}.
\end{itemize}

\nota{El rol \codigo{mpaz\_app} no puede UPDATE/DELETE en las tablas de solo
inserción (\codigo{historial\_cambios}, \codigo{registros\_acceso},
\codigo{administracion\_medicamento}); el script 04 aborta si conserva esos
privilegios. El DDL ya no incluye \codigo{autoriza\_paracetamol}.}

\subsection{Diagramas de datos}

\diagrama{09a-datos-nucleo}{0.45}{Modelo de datos: matrícula, ficha, autorización y prescripción.}
\diagrama{09b-datos-personas}{1.0}{Modelo de datos: personas, documentos y datos sensibles.}
\diagrama{09c-datos-seguridad}{1.0}{Modelo de datos: roles, permisos y auditoría.}
\diagrama{10b-estados-documentos}{0.9}{Estados de la ficha, la autorización y la prescripción.}
\diagrama{10-estados-matricula}{0.52}{Ciclo de la matrícula anual.}
